% =============================================================================
% tema.tex — aspecto de todas las diapositivas salvo la portada.
%
% Mismo estilo que presentacion_referencia/estilo/tema.tex (plantilla PPGE-UCB adaptada): fondo blanco,
% título arriba a la izquierda sin barra, sin cabecera, pie con una franja gruesa (morado) y otra fina
% (celeste), navegación por secciones con un punto por diapositiva y número de diapositiva a la derecha.
% Estilo minimalista: texto en gris oscuro (cInk) y gris (cSrc), un solo acento (cAcc) para títulos y
% cifras clave; los colores de los datos solo en las figuras.
%
% Cambios respecto a la referencia: colores de los estados de la carrera (IDA, FILA, QR, VUELTA), estilos
% TikZ de las figuras del robot, pgfplots para los gráficos, y el dibujo del cierre es el circuito
% (cajas, una señal, fotos con QR y el robot) en lugar de las edificaciones.
% =============================================================================
\usepackage[T1]{fontenc}
\usepackage[utf8]{inputenc}
\usepackage[sfdefault,lining]{FiraSans}
\usepackage[spanish,es-nodecimaldot,es-noshorthands]{babel}
\usepackage{graphicx}
\usepackage{amsmath}
\usepackage{booktabs}
\usepackage{array}
\usepackage{fontawesome5}
\usepackage{tikz}
\usepackage{pgfplots}
\pgfplotsset{compat=1.18}
\usetikzlibrary{positioning,arrows.meta,calc,fit,backgrounds,shapes.misc,shapes.geometric,babel}

% ---- paleta (la de la referencia) -------------------------------------------------
\definecolor{cPub}{HTML}{8A97EE}     % morado
\definecolor{cAcc}{HTML}{4F5FC8}     % morado oscuro (titulos, acentos)
\definecolor{cLight}{HTML}{00BFFF}   % celeste
\definecolor{cCelD}{HTML}{0090C8}    % celeste oscuro (texto sobre blanco)
\definecolor{cEncBg}{HTML}{DEE2FB}   % morado claro
\definecolor{cDecBg}{HTML}{D9F4FF}   % celeste claro
\definecolor{cInk}{HTML}{2E3140}
\definecolor{cOurs}{HTML}{D92626}
\definecolor{cOursD}{HTML}{9E1B1B}
\definecolor{cUni}{HTML}{E08A2E}
\definecolor{cUniD}{HTML}{A85E14}
\definecolor{cSrc}{HTML}{5A5D6E}
\definecolor{cMuted}{HTML}{8A8A96}
\definecolor{cRule}{HTML}{E2E3EA}
\definecolor{cPaper}{HTML}{F5F6FA}
\definecolor{cOk}{HTML}{3C9A5F}
% estados de la carrera (Tarea A), iguales en el circuito, la línea de tiempo y la máquina de estados
% solo dos colores en los diagramas (morado y celeste, con sus tonos claros); los demas son grises
\colorlet{cIda}{cAcc}
\colorlet{cFila}{cCelD}
\colorlet{cQr}{cLight}
\colorlet{cVuelta}{cPub}

% ---- estructura de la plantilla -------------------------------------------------
\useoutertheme[subsection=false]{miniframes}
\useinnertheme{rectangles}
\setbeamertemplate{navigation symbols}{}
\setbeamersize{text margin left=10mm, text margin right=10mm}

\setbeamercolor{normal text}{fg=cInk, bg=white}
\setbeamercolor{background canvas}{bg=white}
\setbeamercolor{structure}{fg=cAcc}
\setbeamercolor{alerted text}{fg=cOurs}
\setbeamercolor{frametitle}{fg=cAcc, bg=white}
\setbeamercolor{footline1}{bg=cPub}
\setbeamercolor{footline2}{bg=cLight}
\setbeamercolor{section in head/foot}{fg=cAcc, bg=white}
\setbeamercolor{mini frame}{fg=cAcc}
\setbeamerfont{frametitle}{size=\Large, series=\mdseries}
\setbeamerfont{section in head/foot}{size=\fontsize{6.5}{7.5}\selectfont}
\setbeamertemplate{itemize item}{\color{cAcc}\raisebox{0.2ex}{\tiny$\blacksquare$}}
\setbeamertemplate{itemize subitem}{\color{cMuted}\raisebox{0.2ex}{\tiny$\blacktriangleright$}}
\setlength{\leftmargini}{1.1em}

\setbeamertemplate{headline}{}

\setbeamertemplate{frametitle}{%
  \vskip 4mm
  {\usebeamerfont{frametitle}\usebeamercolor[fg]{frametitle}\insertframetitle\par}%
  \vskip 0.8mm
}

\makeatletter
\setbeamertemplate{footline}{%
  \begin{beamercolorbox}[wd=\paperwidth, ht=2.2mm, dp=0pt]{footline1}\end{beamercolorbox}%
  \begin{beamercolorbox}[wd=\paperwidth, ht=0.6mm, dp=0pt]{footline2}\end{beamercolorbox}%
  \begin{beamercolorbox}[wd=\paperwidth]{section in head/foot}%
    \vskip 2pt
    \hbox to \paperwidth{\insertnavigation{\dimexpr\paperwidth-14mm\relax}\hfill
      \raisebox{0.4mm}{\fontsize{6}{7}\selectfont\color{cSrc}\insertframenumber\,/\,\inserttotalframenumber}\hskip 5mm}%
    \vskip 3pt
  \end{beamercolorbox}}
\makeatother

% ---- piezas reutilizables ------------------------------------------------------------------
% cifra grande con leyenda: \cifra{numero}{leyenda}
\newcommand{\cifra}[2]{{\color{cAcc}\fontsize{17}{19}\selectfont #1}\par\vspace{0.3mm}%
  {\scriptsize\color{cSrc}#2\par}}
% texto secundario
\newcommand{\nota}[1]{{\color{cSrc}#1}}
% pie de figura (un párrafo chico) y encabezado de una columna
\newcommand{\etq}[1]{{\fontsize{6.3}{7.6}\selectfont\color{cSrc}#1\par}}
\newcommand{\cab}[1]{{\footnotesize #1\par}\vspace{0.8mm}}

\tikzset{
  % --- diagramas de bloques
  caja/.style={draw=cInk!45, line width=0.5pt, fill=white, rounded corners=2pt, align=center,
    inner sep=1.8pt, font=\fontsize{6.8}{8}\selectfont, text=cInk, execute at begin node={\hyphenpenalty=10000}},
  grupo/.style={draw=cInk!25, line width=0.5pt, fill=cPaper, rounded corners=3pt},
  ar/.style={-{Stealth[length=1.6mm]}, line width=0.6pt, draw=cInk!55},
  arAcc/.style={-{Stealth[length=1.6mm]}, line width=0.8pt, draw=cAcc},
  lbl/.style={font=\fontsize{6}{7}\selectfont, text=cSrc, align=center},
  img/.style={inner sep=0pt, draw=cRule, line width=0.4pt},
  % --- estados de la carrera
  est/.style={draw=#1, line width=0.8pt, fill=#1!10, rounded corners=6pt, font=\scriptsize\bfseries, text=#1!70!black,
    minimum height=6mm, inner xsep=2.5pt, align=center},
  % --- circuito visto desde arriba (figuras/gen/circuito.tex, coordenadas en m)
  pared/.style={draw=cInk!70, line width=1.3pt, line cap=round},
  senal/.style={circle, fill=white, draw=cAcc, line width=0.8pt, inner sep=0pt, minimum size=2.4mm},
  foto/.style={rectangle, fill=cMuted, draw=white, line width=0.3pt, inner sep=0pt, minimum size=1.9mm},
  buscada/.style={rectangle, fill=cAcc, draw=white, line width=0.3pt, inner sep=0pt, minimum size=2.4mm},
  puntoA/.style={diamond, fill=cInk, draw=white, line width=0.4pt, inner sep=0pt, minimum size=3mm},
  % la bandera a cuadros del fondo (la meta): un tablero de 2 x 2
  bandera/.style={rectangle, fill=white, draw=cInk, line width=0.5pt, inner sep=0pt, minimum size=3mm,
    path picture={\fill[cInk] (path picture bounding box.south west) rectangle (path picture bounding box.center)
      (path picture bounding box.center) rectangle (path picture bounding box.north east);}},
  tIda/.style={draw=cIda, line width=1pt, line join=round},
  tFila/.style={draw=cFila, line width=0.8pt, line join=round},
  tQr/.style={draw=cQr, line width=1.3pt, line join=round},
  tVuelta/.style={draw=cVuelta, line width=1pt, line join=round, dash pattern=on 2.2pt off 1.2pt},
  % --- ventana dinamica (figuras/gen/arcos.tex, marco del robot, en m)
  arcoOk/.style={draw=cAcc!45, line width=0.35pt},
  arcoMal/.style={draw=cMuted!80, line width=0.35pt},
  arcoElegido/.style={draw=cAcc, line width=1.6pt, -{Stealth[length=1.6mm]}},
}

% ---- pgfplots: ejes sobrios, en la fuente y los colores de la presentación
\pgfplotsset{
  sobrio/.style={
    axis line style={draw=cInk!45}, tick style={draw=cInk!45},
    tick label style={font=\fontsize{6}{7}\selectfont, text=cSrc},
    label style={font=\fontsize{6.5}{7.5}\selectfont, text=cSrc},
    legend style={font=\fontsize{6}{7}\selectfont, draw=none, fill=none, text=cInk},
    grid style={draw=cRule}, major grid style={draw=cRule},
    every axis plot/.append style={line width=0.9pt},
  },
}

% ---- dibujo lineal del cierre: el circuito (cajas, una señal, fotos con QR y el robot) ----------------
% unidades en mm, suelo en y = 0
\newcommand{\skCaja}[3]{% x, ancho, alto
  \draw[skL, line width=0.45pt, fill=white] (#1,0) rectangle ++(#2,#3);
  \draw[skL, line width=0.3pt] (#1,{#3*0.82}) -- ++(#2,0);
  \draw[skL, line width=0.3pt] ({#1+#2/2},{#3*0.82}) -- ++(0,{#3*0.18});}
\newcommand{\skFoto}[2]{% x, y: foto con su QR debajo, pegada a una caja
  \draw[cAcc, line width=0.4pt, fill=cEncBg] (#1,#2) rectangle ++(4,5);
  \draw[cAcc, line width=0.4pt] ({#1+2},{#2+3.3}) circle (0.9);
  \draw[cAcc, line width=0.4pt] ({#1+0.6},{#2+0.4}) .. controls ({#1+1},{#2+2}) and ({#1+3},{#2+2}) .. ({#1+3.4},{#2+0.4});
  \draw[skL, line width=0.4pt, fill=white] ({#1+0.5},{#2-4.2}) rectangle ++(3,3);
  \fill[skL] ({#1+0.8},{#2-3.9}) rectangle ++(0.9,0.9);
  \fill[skL] ({#1+2.3},{#2-3.9}) rectangle ++(0.9,0.9);
  \fill[skL] ({#1+0.8},{#2-2.4}) rectangle ++(0.9,0.9);
  \fill[skL] ({#1+2.4},{#2-2.3}) rectangle ++(0.5,0.5);}
\newcommand{\skRobot}[1]{% x: TurtleBot 4 Lite de costado
  \draw[skL, line width=0.5pt, fill=white, rounded corners=1.5pt] (#1,1) rectangle ++(14,5);
  \draw[skL, line width=0.5pt, fill=white] ({#1+4},6) rectangle ++(6,1.6);
  \draw[skL, line width=0.5pt, fill=white] ({#1+5.5},7.6) rectangle ++(3,1.4);
  \draw[cAcc, line width=0.5pt, fill=cEncBg] ({#1+12.2},3.2) rectangle ++(1.8,1.6);
  \draw[skL, line width=0.5pt, fill=white] ({#1+3},1.2) circle (1.2);
  \draw[skL, line width=0.5pt, fill=white] ({#1+11},1.2) circle (1.2);
  \draw[cLight, line width=0.6pt, dash pattern=on 1pt off 0.8pt] ({#1+14.3},4) -- ++(5,1.6);
  \draw[cLight, line width=0.6pt, dash pattern=on 1pt off 0.8pt] ({#1+14.3},4) -- ++(5,-1.2);}
\newcommand{\circuitoLinea}{%
  \begin{tikzpicture}[x=1mm, y=1mm]
    \colorlet{skL}{cInk!45}
    \draw[skL, line width=0.5pt] (-2,0) -- (128,0);
    \skCaja{0}{14}{11}
    \skCaja{14}{12}{8}
    \draw[cAcc, line width=0.5pt, fill=white] (18,13) circle (3.2);
    \draw[cAcc, line width=0.7pt, -{Stealth[length=1.3mm]}] (16.6,11.6) -- (16.6,13.6) -- (19.6,13.6);
    \draw[cInk!30, line width=0.6pt, dash pattern=on 1.5pt off 1.5pt] (34,0.5) .. controls (50,0.5) and (52,0.5) .. (62,0.5);
    \skRobot{40}
    \skCaja{68}{12}{14}
    \skCaja{80}{14}{14}
    \skCaja{94}{12}{14}
    \skCaja{106}{16}{14}
    \skFoto{72}{6.5}
    \skFoto{85}{6.5}
    \skFoto{97}{6.5}
    \skFoto{112}{6.5}
    \draw[cCelD, line width=0.8pt, rounded corners=0.8pt] (84.2,1.6) rectangle (89.8,12.3);
  \end{tikzpicture}}
